\documentclass{article}

\usepackage{amsmath,amssymb}
\usepackage{xurl}
\usepackage[hidelinks]{hyperref}
\setlength{\emergencystretch}{2em}

\input{command}

\title{Zero Multiplicities in Wolf's Two-Pure-State Spectrum Shorthand}
\author{}
\date{2026-08-08}

\begin{document}
\maketitle
\gapnote{local-correction}{resolved}

\section{The source display}

Let $P,Q\in M_d(\C)$ be Hermitian rank-one projections. In the
spectrum-preserver argument following Wigner's theorem, Wolf writes the
spectrum of $P+Q$ as
\begin{align}
    1-\sqrt{\tr(PQ)},
    \qquad
    1+\sqrt{\tr(PQ)}.
    \label{eq:wolf_ch01_two_pure_state_zero_multiplicities_source_two_values}
\end{align}
This shorthand appears in Chapter~1, lines~289--296 of
Ref.~\cite{Wolf2012Quantum}; the local source is
\path{Notes/WolfNoteTexSource/ch01_deconstructing_quantum.tex}.

\section{The suppressed eigenvalues}

The range of $P+Q$ lies in the span of the one-dimensional ranges of $P$ and
$Q$. Hence $P+Q$ vanishes on a subspace of dimension at least $d-2$. For
$d\geq2$, its characteristic polynomial, with algebraic multiplicities, is
\begin{align}
    \chi_{P+Q}(X)
        &=X^{d-2}((X-1)^2-\tr(PQ)).
    \label{eq:wolf_ch01_two_pure_state_zero_multiplicities_characteristic_polynomial}
\end{align}
The two values in
(\ref{eq:wolf_ch01_two_pure_state_zero_multiplicities_source_two_values}) are
the roots of the quadratic factor in
(\ref{eq:wolf_ch01_two_pure_state_zero_multiplicities_characteristic_polynomial}).
The factor $X^{d-2}$ contributes at least $d-2$ zero eigenvalues.

If $P=Q$, then $\tr(PQ)=1$, so
\begin{align}
    (X-1)^2-\tr(PQ)
        &=X(X-2).
    \label{eq:wolf_ch01_two_pure_state_zero_multiplicities_equal_projection_quadratic}
\end{align}
The quadratic factor then contributes one additional zero, and the total zero
multiplicity is $d-1$.

In dimension one, every normalized rank-one projection is the identity, so the
separate formula is
\begin{align}
    \chi_{P+Q}(X)
        &=X-2.
    \label{eq:wolf_ch01_two_pure_state_zero_multiplicities_dimension_one_charpoly}
\end{align}
In dimension zero, projective pure states do not exist.

\section{Formal correction}

The formalization proves the multiplicity-aware identity
(\ref{eq:wolf_ch01_two_pure_state_zero_multiplicities_characteristic_polynomial})
through a $d\times2$ and $2\times d$ factorization. It then recovers
$\tr(PQ)$ from equality of two such characteristic polynomials.\footnote{Formal
declarations:
\leanid{Projectivization.charpoly_add_pureStateMatrix_of_two_le} and
\leanid{Projectivization.trace_pureStateMatrix_mul_eq_of_charpoly_add_eq}
in \doclink{QICLean/Channel/Wigner/TwoPureStateCharpoly}.}
The low-dimensional cases are stated explicitly in the same file. This
correction is tracked by
\href{https://github.com/LionSR/TNLean/issues/5683}{TNLean issue \#5683}.

\section{Classification}

\emph{Local correction.} The issue is a spectrum shorthand, not a failure of
Wolf's intended spectrum-preserver argument \cite{Wolf2012Quantum}. The two
roots in
(\ref{eq:wolf_ch01_two_pure_state_zero_multiplicities_source_two_values})
retain the transition-probability information used in that argument. The
formal theorem uses characteristic polynomials because they record every zero
eigenvalue and every algebraic multiplicity without ambiguity.

\bibliographystyle{alpha}
\bibliography{references}

\end{document}